\documentclass[12pt]{article}
\usepackage{amsmath,amssymb,amsfonts,amsthm}
\usepackage[margin=1in]{geometry}

\begin{document}

\begin{center}
{\Large\bfseries Chapter 9, Problem 2}\\[6pt]
Byron \& Fuller, \textit{Mathematics of Classical and Quantum Physics}
\end{center}

\bigskip

\textbf{Problem.} Show that if we modify $k(s,t)$ on a set of measure zero in the plane, then the resulting kernel $\tilde{k}(s,t)$ has essentially the same eigenvectors and eigenvalues as $k(s,t)$, and the solution to the integral equation
\[
f(x) = g(x) + \lambda \int k(x,t)\,f(t)\,dt
\]
is essentially the same as the solution to the corresponding equation with $\tilde{k}$. What is meant by ``essentially'' in the above statement?

\bigskip
\textbf{Solution.}

\textbf{Meaning of ``essentially'':} Two functions that differ only on a set of measure zero are equal as elements of $L_2$. That is, they are in the same equivalence class. ``Essentially the same'' means ``equal almost everywhere,'' i.e., equal as elements of $L_2$.

\textbf{Proof:} Let $\tilde{k}(s,t) = k(s,t)$ except on a set $E \subset [a,b] \times [a,b]$ of (two-dimensional Lebesgue) measure zero. Define the difference kernel $\delta k(s,t) = \tilde{k}(s,t) - k(s,t)$. Then $\delta k(s,t) = 0$ for all $(s,t) \notin E$, i.e., $\delta k = 0$ almost everywhere on $[a,b]^2$.

Consider the integral operator $\delta K$ with kernel $\delta k$:
\[
(\delta K f)(x) = \int_a^b \delta k(x,t)\,f(t)\,dt.
\]

For any $f \in L_2[a,b]$, we have by the Cauchy--Schwarz inequality:
\[
|(\delta K f)(x)|^2 \le \int_a^b |\delta k(x,t)|^2\,dt \cdot \int_a^b |f(t)|^2\,dt = \|f\|^2 \int_a^b |\delta k(x,t)|^2\,dt.
\]

Therefore:
\[
\|\delta K f\|^2 = \int_a^b |(\delta K f)(x)|^2\,dx \le \|f\|^2 \int_a^b \int_a^b |\delta k(x,t)|^2\,dx\,dt = \|f\|^2 \cdot \|\delta k\|_{L_2}^2.
\]

Since $\delta k = 0$ a.e.\ on $[a,b]^2$, we have $\|\delta k\|_{L_2}^2 = 0$. Therefore $\|\delta K f\| = 0$ for all $f$, which means $\delta K = 0$ as an operator on $L_2$.

Thus $\tilde{K} = K + \delta K = K$ as operators on $L_2$. Since the two integral operators are identical as operators on $L_2$:

\begin{enumerate}
\item They have the same eigenvalues.
\item Their eigenvectors are the same (as elements of $L_2$, i.e., equal a.e.).
\item The solutions to $f = g + \lambda K f$ and $f = g + \lambda \tilde{K} f$ are the same as elements of $L_2$.
\end{enumerate}

In summary, modifying the kernel on a set of measure zero does not change the operator at all in $L_2$, so all spectral properties and solutions are preserved. The word ``essentially'' refers to equality in the $L_2$ sense (equality almost everywhere). \hfill $\blacksquare$

\end{document}
